%========================================
% MatinBook - Stage 01: Basic Test
% Test: Core System + Colors + Fonts + Math
% Purpose: Validate basic infrastructure works
%========================================

% Add parent directory to search path
\makeatletter
\def\input@path{{../}}
\makeatother

\documentclass{matinbook}

\title{آزمون زیرساخت پایه}
\author{تیم توسعه MatinBook}
\date{تابستان ۱۴۰۵}

\begin{document}

%----------------------------------------
% Title Page
%----------------------------------------
\maketitle

%----------------------------------------
% Chapter 1: Color System
%----------------------------------------
\chapter{سیستم رنگ‌ها}

\section{پالت رنگی اصلی}

رنگ‌های اصلی \lr{MatinBook} برای تفکیک بصری انواع محتوا طراحی شده‌اند.
هر رنگ خانواده‌ای از سه طیف روشن، اصلی و تیره دارد.

\subsection{رنگ‌های آبی (قضایا)}

خانواده آبی برای قضایا، لم‌ها و نتایج:

\begin{itemize}
    \item \textcolor{matinblue}{\textbf{آبی اصلی}} --- \texttt{matinblue} (RGB: 41, 128, 185)
    \item \textcolor{matindarkblue}{\textbf{آبی تیره}} --- \texttt{matindarkblue} (RGB: 21, 67, 96)
    \item \colorbox{matinlightblue}{\textbf{آبی روشن}} --- \texttt{matinlightblue} (RGB: 214, 234, 248)
\end{itemize}

\subsection{رنگ‌های سبز (تعاریف)}

خانواده سبز برای تعاریف و مفاهیم:

\begin{itemize}
    \item \textcolor{matingreen}{\textbf{سبز اصلی}} --- \texttt{matingreen} (RGB: 39, 174, 96)
    \item \textcolor{matindarkgreen}{\textbf{سبز تیره}} --- \texttt{matindarkgreen} (RGB: 20, 90, 50)
    \item \colorbox{matinlightgreen}{\textbf{سبز روشن}} --- \texttt{matinlightgreen} (RGB: 215, 245, 227)
\end{itemize}

\subsection{رنگ‌های نارنجی (مثال‌ها)}

خانواده نارنجی برای مثال‌ها:

\begin{itemize}
    \item \textcolor{matinorange}{\textbf{نارنجی اصلی}} --- \texttt{matinorange} (RGB: 230, 126, 34)
    \item \textcolor{matindarkorange}{\textbf{نارنجی تیره}} --- \texttt{matindarkorange} (RGB: 120, 66, 18)
    \item \colorbox{matinlightorange}{\textbf{نارنجی روشن}} --- \texttt{matinlightorange} (RGB: 252, 235, 213)
\end{itemize}

\subsection{رنگ‌های قرمز (نکات و هشدارها)}

خانواده قرمز برای نکات مهم:

\begin{itemize}
    \item \textcolor{matinred}{\textbf{قرمز اصلی}} --- \texttt{matinred} (RGB: 231, 76, 60)
    \item \textcolor{matindarkred}{\textbf{قرمز تیره}} --- \texttt{matindarkred} (RGB: 120, 40, 31)
    \item \colorbox{matinlightred}{\textbf{قرمز روشن}} --- \texttt{matinlightred} (RGB: 250, 215, 210)
\end{itemize}

\subsection{رنگ‌های بنفش (تمرین‌ها)}

خانواده بنفش برای تمرین‌ها:

\begin{itemize}
    \item \textcolor{matinpurple}{\textbf{بنفش اصلی}} --- \texttt{matinpurple} (RGB: 142, 68, 173)
    \item \textcolor{matindarkpurple}{\textbf{بنفش تیره}} --- \texttt{matindarkpurple} (RGB: 74, 35, 90)
    \item \colorbox{matinlightpurple}{\textbf{بنفش روشن}} --- \texttt{matinlightpurple} (RGB: 235, 222, 245)
\end{itemize}

\subsection{رنگ‌های خنثی}

\begin{itemize}
    \item \textcolor{matindarkgray}{\textbf{خاکستری تیره}} --- متن اصلی
    \item \textcolor{matinmediumgray}{\textbf{خاکستری متوسط}} --- متن فرعی
    \item \colorbox{matinlightgray}{\textbf{خاکستری روشن}} --- پس‌زمینه
    \item \colorbox{matinbordergray}{\textbf{خاکستری مرزی}} --- حاشیه‌ها
\end{itemize}

\section{نمودار رنگی}

\begin{center}
\begin{latin}
\begin{tikzpicture}
    % Color swatches
    \foreach \color [count=\i] in {matinblue, matingreen, matinorange, matinred, matinpurple} {
        \fill[\color] (\i*2.5-1.5, 2) rectangle (\i*2.5-0.5, 3);
        \fill[\color!50] (\i*2.5-1.5, 0.5) rectangle (\i*2.5-0.5, 1.5);
        \fill[\color!20] (\i*2.5-1.5, -1) rectangle (\i*2.5-0.5, 0);
    }
\end{tikzpicture}
\end{latin}
\end{center}

%----------------------------------------
% Chapter 2: Font System
%----------------------------------------
\chapter{سیستم فونت‌ها}

\section{فونت فارسی}

این یک پاراگراف طولانی فارسی است برای تست فونت پیش‌فرض \lr{XB Niloofar}.
این فونت باید خوانایی خوبی داشته باشد و تمام حروف فارسی را
به درستی نمایش دهد. حروف خاص فارسی مانند «گ»، «چ»، «پ»، «ژ»
باید به صورت صحیح نمایش داده شوند.

اعداد فارسی: ۰ ۱ ۲ ۳ ۴ ۵ ۶ ۷ ۸ ۹

اعداد انگلیسی در متن فارسی: 0 1 2 3 4 5 6 7 8 9

نویسه‌های ویژه فارسی: «نقل قول»، (پرانتز)، [کروشه]، \{آکولاد\}

\section{فونت انگلیسی (لاتین)}

\begin{latin}
This is an English paragraph to test the Latin font Times New Roman.
It should render professionally with proper kerning and ligatures.

The quick brown fox jumps over the lazy dog.
Pack my box with five dozen liquor jugs.

Special ligatures: fi fl ff ffi ffl --- -- ``quotes''

Mathematical symbols in text: \( \alpha, \beta, \gamma, \delta \)
\end{latin}

\section{ترکیب فارسی و انگلیسی}

این یک متن فارسی است که شامل کلمات انگلیسی مانند \lr{computer}،
\lr{software}، \lr{algorithm}، \lr{function} و \lr{variable} می‌باشد.
همچنین اعداد لاتین مانند \lr{3.14159} و \lr{42} در متن فارسی قرار گرفته‌اند.

نام‌های خاص خارجی: \lr{Donald Knuth}، \lr{Leslie Lamport}، \lr{Alan Turing}

\section{فونت ریاضی}

آزمایش فونت ریاضی \lr{XITS Math} با معادلات پایه:

معادله درون خطی: $E = mc^2$، قضیه فیثاغورث: $a^2 + b^2 = c^2$

معادله نمایشی:
\begin{equation}
    \int_{0}^{\infty} \frac{x^3}{e^x - 1} \,dx = \frac{\pi^4}{15}
    \label{eq:stefan-boltzmann}
\end{equation}

حروف یونانی: $\alpha, \beta, \gamma, \Gamma, \delta, \Delta, \epsilon, \varepsilon,
\zeta, \eta, \theta, \Theta, \lambda, \Lambda, \mu, \pi, \Pi, \sigma, \Sigma, \tau,
\phi, \Phi, \omega, \Omega$

عملگرهای بزرگ:
\[
    \sum_{n=1}^{\infty} \frac{1}{n^2} = \frac{\pi^2}{6}, \qquad
    \prod_{k=1}^{n} k = n!, \qquad
    \lim_{x \to \infty} \left(1 + \frac{1}{x}\right)^x = e
\]

%----------------------------------------
% Chapter 3: Color Boxes Test
%----------------------------------------
\chapter{جعبه‌های رنگی}

\section{نمایش رنگ‌ها در جعبه‌های اطلاعاتی}

\begin{theorem}[قضیه تست]
    این یک قضیه آزمایشی با رنگ آبی است.
    تمام عناوین باید خوانا و متمایز باشند.
\end{theorem}

\begin{definition}[تعریف تست]
    این یک تعریف آزمایشی با رنگ سبز است.
    رنگ‌ها باید در حالت چاپ سیاه و سفید نیز قابل تشخیص باشند.
\end{definition}

\begin{example}[مثال تست]
    این یک مثال آزمایشی با رنگ نارنجی است.
    مثال‌ها باید توجه خواننده را به کاربرد مفاهیم جلب کنند.
\end{example}

\begin{remark}[نکته تست]
    این یک نکته آزمایشی با رنگ قرمز است.
    نکات مهم باید برجسته و قابل توجه باشند.
\end{remark}

\begin{exercise}[تمرین تست]
    این یک تمرین آزمایشی با رنگ بنفش است.
    تمرین‌ها باید خواننده را به تفکر وادارند.
\end{exercise}

%----------------------------------------
% Chapter 4: Typography
%----------------------------------------
\chapter{تایپوگرافی}

\section{فاصله خطوط و ترازبندی}

این یک پاراگراف طولانی فارسی است که برای تست تایپوگرافی نوشته شده است.
هدف ما این است که ببینیم آیا خطوط به خوبی تراز می‌شوند و فاصله بین
کلمات مناسب است یا خیر. تایپوگرافی خوب باعث افزایش خوانایی متن
و کاهش خستگی چشم خواننده می‌شود. ما در \lr{MatinBook} از تنظیمات
پیشرفته \lr{microtype} برای بهبود تایپوگرافی استفاده می‌کنیم.

فاصله بین خطوط باید به اندازه‌ای باشد که خواننده به راحتی بتواند
خط بعدی را پیدا کند، اما نه آنقدر زیاد که متن از هم گسیخته به نظر برسد.
همچنین ترازبندی دوطرفه (\lr{justified}) باید به گونه‌ای باشد که
فضای خالی بیش از حد بین کلمات ایجاد نشود.

\section{جلوگیری از خطوط آویزان}

در حروف‌چینی حرفه‌ای، باید از ایجاد خطوط آویزان (\lr{widow}) و
یتیم (\lr{orphan}) جلوگیری شود. خط آویزان به آخرین خط یک پاراگراف
گفته می‌شود که به تنهایی در بالای صفحه بعد قرار گیرد. خط یتیم
به اولین خط یک پاراگراف گفته می‌شود که در انتهای صفحه باقی بماند.

\lr{MatinBook} با تنظیم \texttt{\textbackslash widowpenalty=10000}
و \texttt{\textbackslash clubpenalty=10000} از این مشکلات جلوگیری می‌کند.

%----------------------------------------
% System Check Summary
%----------------------------------------
\chapter{خلاصه نتایج}

\section{وضعیت ماژول‌ها}

\begin{itemize}
    \item \textcolor{matingreen}{\textbf{سیستم رنگ:}} ۱۵ رنگ در ۵ خانواده --- ✅
    \item \textcolor{matingreen}{\textbf{فونت فارسی:}} \lr{XB Niloofar} با مقیاس ۱.۲ --- ✅
    \item \textcolor{matingreen}{\textbf{فونت لاتین:}} \lr{Times New Roman} --- ✅
    \item \textcolor{matingreen}{\textbf{فونت ریاضی:}} \lr{XITS Math} --- ✅
    \item \textcolor{matingreen}{\textbf{فونت کد:}} \lr{JetBrains Mono} --- ✅
    \item \textcolor{matingreen}{\textbf{جعبه‌های رنگی:}} ۱۰ نوع مختلف --- ✅
    \item \textcolor{matinggreen}{\textbf{تایپوگرافی:}} \lr{microtype} فعال --- ✅
    \item \textcolor{matingreen}{\textbf{ارجاع هوشمند:}} \lr{cleveref} فارسی --- ✅
\end{itemize}

\section{نتیجه}

\textbf{ماژول‌های پایه \lr{MatinBook} نسخه ۱ با موفقیت بارگذاری و آزمایش شدند.}
تمام سیستم‌های رنگی، فونت‌ها و جعبه‌ها به درستی کار می‌کنند. 🎉

\end{document}